% !TeX root = ../../Book4.tex
% 纲要标题：### 13.4 典型计算、Morita 不变性与 HKR
% 纲要标签：b4:sec:1204
% 纲要位置：计划纲要.md，第 3549 行。

\section{典型计算、Morita 不变性与 HKR}
\label{b4:sec:1204}

矩阵代数、光滑交换代数及一个平方零代数展示三种不同的 Hochschild 计算。前两种分别利用模范畴等价和对角的 Koszul 消解，最后一种直接检验形变障碍。

% 纲要内容（仅作写作计划，尚非教材正文）：
%
% * 矩阵代数与 Morita 不变性
% * 光滑交换代数的比较方向
% * 不把一阶可形变误判为一定可延拓至所有阶
%
% ---
%

\subsection{矩阵代数与 Morita 不变性}
\label{b4:subsec:1204:01}

矩阵代数的乘法看起来比标量代数复杂，但模范畴之间的等价会保留相应的同调信息。下面用矩阵单位写出这一等价，并从零次交换子商的比较得到 Hochschild 同调的 Morita 不变性\footnote{\OriginalHanName{森田紀一}（假名：もりた きいち，罗马音：Kiiti Morita，1915-1995），日本数学家，建立以模范畴等价为核心的森田理论。}。

\begin{theorem}\label{b4:thm:hh-matrix-morita}
设 \(k\) 为域，\(A\) 为结合含幺 \(k\) 代数，\(r\ge1\)，\(B=M_r(A)\)。对任意 \(A\) 双模 \(N\)，令 \(M_r(N)\) 带有左右矩阵乘法作用。则对所有 \(n\ge0\)，有自然同构
\[
\operatorname{HH}_n(B,M_r(N))\cong\operatorname{HH}_n(A,N),\qquad
\operatorname{HH}^n(B,M_r(N))\cong\operatorname{HH}^n(A,N).
\]
\end{theorem}
\begin{proof}
记 \(E_{ij}\) 为矩阵单位，\(e=E_{11}\)，并识别 \(eBe=A\)。对 \(M=M_r(N)\)，\(E_{1i}mE_{j1}\) 正是将 \(m\) 的第 \((i,j)\) 个分量放进 \((1,1)\) 角部，所以该角部足以恢复整个矩阵双模。对一般 \(B\) 双模，仍可用相同的矩阵单位来恢复各分量。双模函子
\[
F(M)=eMe,\qquad G(N)=M_r(N)
\]
互为等价。具体地，\(FG(N)=N\)，而
\[
M\longrightarrow M_r(eMe),\qquad
m\longmapsto(E_{1i}mE_{j1})_{i,j}
\]
的逆将 \((n_{ij})\) 送到 \(\sum_{i,j}E_{i1}n_{ij}E_{1j}\)。矩阵单位等式 \(E_{ij}E_{uv}=\delta_{ju}E_{iv}\) 验证两次复合为恒等，且这些映射保持双模作用。两个函子均正合：若 \(f:M\twoheadrightarrow M'\) 为双模满射，给定 \(em'e\)，先取 \(f(m)=m'\)，再取角部 \(eme\)，就有 \(f(eme)=em'e\)；核则满足 \(\ker F(f)=e(\ker f)e\)。而 \(M_r(-)\) 是有限直和，也保持正合。它们还保持投射对象，因为一个提升问题经逆等价转移后仍是满射的提升问题。

对 \(B\) 取投射双模消解 \(P_\bullet\to B\)。施加 \(F\) 后得到 \(A\) 的投射双模消解，而且等价的满忠实性给出复形同构
\[
\operatorname{Hom}_{B^e}(P_\bullet,M)
\cong\operatorname{Hom}_{A^e}(F(P_\bullet),F(M)).
\]
取上同调便证明上同调断言。

上同调刚才使用的是代数本身的消解；同调则可以借\cref{b4:thm:tor-balanced}改用系数双模的消解。这样便能逐项使用双模范畴等价。为此先考察右正合函子 \(Q_B(M)=M/[B,M]\)。它与张量函子的联系由以下互逆映射给出：
\[
\begin{aligned}
M^{\mathrm r}\otimes_{B^e}B&\longrightarrow Q_B(M),
&m\otimes b&\longmapsto[bm],\\
Q_B(M)&\longrightarrow M^{\mathrm r}\otimes_{B^e}B,
&[m]&\longmapsto m\otimes1.
\end{aligned}
\]
第一式满足平衡关系：对 \(a,b,c\in B\)，\((cma)\otimes b=m\otimes abc\)，而交换子商中 \([bcma]=[abcm]\)。第二式消去交换子，因为
\[
\begin{aligned}
bm\otimes1
&=m\cdot(1\otimes b^{\mathrm{op}})\otimes1
=m\otimes b,\\
mb\otimes1
&=m\cdot(b\otimes1)\otimes1
=m\otimes b.
\end{aligned}
\]
此外 \(bm\otimes1=m\otimes b\) 验证两次复合为恒等。因此 \(Q_B(M)\cong M^{\mathrm r}\otimes_{B^e}B\)。从双模到右 \(B^e\) 模的函子 \(M\mapsto M^{\mathrm r}\) 是正合等价，保持投射对象，所以在系数双模 \(M\) 的投射消解上计算，并用\cref{b4:thm:tor-balanced}，得到
\[
L_nQ_B(M)
\cong\operatorname{Tor}^{B^e}_n(M^{\mathrm r},B)
=\operatorname{HH}_n(B,M).
\]

现在比较两个交换子商。矩阵单位给出自然同构
\[
Q_B(M)\longrightarrow Q_A(eMe),\qquad
[m]\longmapsto\left[\sum_iE_{1i}mE_{i1}\right].
\]
其良定义可直接检验。将 \(1=\sum_jE_{j1}E_{1j}\) 插入 \(b\) 与 \(m\) 之间，得到
\[
E_{1i}bmE_{i1}=\sum_j(E_{1i}bE_{j1})(E_{1j}mE_{i1}).
\]
在目标交换子商中交换这两个因子的先后，再交换求和指标 \(i,j\)，就得到 \(\sum_iE_{1i}mbE_{i1}\)。逆映射由角部的包含 \(eMe\subseteq M\) 给出。在 \(Q_B(M)\) 中，关系
\[
E_{ii}m-E_{1i}mE_{i1}=[E_{i1},E_{1i}m]
\]
说明逆向复合是恒等；正向复合在 \(eMe\) 上也显然为恒等。

所得零次比较就是 \(Q_B\cong Q_A\circ F\)。\(F\) 正合且保持投射，将 \(M\) 的双模投射消解送到 \(F(M)\) 的投射消解；逐项应用零次自然同构，再取同调，便有
\[
\operatorname{HH}_n(B,M)
\cong L_nQ_B(M)
\cong L_nQ_A(F(M))
\cong\operatorname{HH}_n(A,F(M)).
\]
取 \(M=M_r(N)\) 即得结论。
\end{proof}

在 \(M=M_r(N)\) 时，刚才的零次比较把矩阵 \((m_{ij})\) 的类送到 \([\sum_i m_{ii}]\)，即由矩阵迹给出；矩阵单位 \(E_{1i}mE_{i1}\) 将每个对角分量都移到同一角部。这里的迹取值于 \(N/[A,N]\)，其对循环换序的不变性来自取商，未要求 \(A\) 交换。

\ProseParagraphBreak
一般的 \(k\) 线性 Morita 等价由互逆的张量双模实现，同样保持相应系数的 Hochschild 群。矩阵情形的具体等价与\cite[Proposition 9.5.2, Corollary 9.5.3]{Weibel1994HomologicalAlgebra}一致；下面说明一般情形的零次比较为何也能导出。

\begin{proofsketch}
设 \({}_AP_B\)、\({}_BQ_A\) 实现第二卷定理16.3.3及推论16.3.4的 Morita 等价，取相容的双模同构
\(P\otimes_BQ\cong A\)、\(Q\otimes_AP\cong B\)，分别把配对写作 \(\langle p,q\rangle_A\) 与 \(\langle q,p\rangle_B\)。双模函子
\[
 F(N)=Q\otimes_AN\otimes_AP
\]
及其逆是正合等价，送 \(A\) 到 \(B\)，并保持投射对象。这里选择从 \(A\) 双模到 \(B\) 双模的方向；它与刚才从矩阵双模取角部的方向相反，而在矩阵情形对应 \(N\mapsto M_r(N)\)。因此对 \(A\) 的双模投射消解施加 \(F\)，再用等价的 Hom 双射，得到
\(\operatorname{HH}^n(A,N)\cong\operatorname{HH}^n(B,F(N))\)。

同调端的零次比较为
\[
 F(N)/[B,F(N)]\longrightarrow N/[A,N],\qquad
 [q\otimes n\otimes p]\longmapsto
 [n\langle p,q\rangle_A].
\]
张量平衡关系在右边分别变成等式或交换子关系；\(B\) 交换子关系由
\(\langle pb,q\rangle_A=\langle p,bq\rangle_A\) 消去。取有限表达
\(1_A=\sum_i\langle p_i,q_i\rangle_A\)。这是单位在 \(P\otimes_BQ\) 中的原像，而普通张量积的每个元素都是有限个纯张量之和，所以这里的有限性不要求 \(A\) 有限维。逆映射送 \([n]\) 到
\(\sum_i[q_i\otimes n\otimes p_i]\)。一向复合直接为恒等，另一向使用 Morita 配对的相容式
\[
 \langle p,q\rangle_Ap_i=p\langle q,p_i\rangle_B,
 \qquad
 \langle q,p_i\rangle_Bq_i=q\langle p_i,q_i\rangle_A,
\]
以及交换子商中允许把右侧的 \(B\) 因子移到左侧，便恢复 \([q\otimes n\otimes p]\)。这也使逆与所选单位表达无关。

上述张量识别对任意结合代数都成立，因而
\[
L_nQ_A(N)\cong\operatorname{HH}_n(A,N),
\qquad
L_nQ_B(F(N))\cong\operatorname{HH}_n(B,F(N)).
\]
在 \(N\) 的双模投射消解上逐项应用 \(Q_B\circ F\cong Q_A\)，再利用 \(F\) 正合且保持投射，就得到两端 Hochschild 同调的同构。这里省略的是把等价的单位、余单位写成上述相容配对的标准 Morita 构造。完整同调比较也见\cite[Theorem 9.5.6, pp. 328--329]{Weibel1994HomologicalAlgebra}，该处用双复形比较实现相同的系数变换。
\end{proofsketch}

\begin{example}\label{b4:exm:matrix-hochschild-vanishing}
对任意域 \(k\) 和 \(r\ge1\)，\(k^e=k\)，所以
\[
\begin{gathered}
\operatorname{HH}_0(M_r(k))=k,\qquad
\operatorname{HH}^0(M_r(k))=k,\\
\operatorname{HH}_n(M_r(k))=\operatorname{HH}^n(M_r(k))=0\quad(n>0).
\end{gathered}
\]
零次同调的同构是普通迹在交换子商上诱导的映射，其逆送 \(c\) 到 \([cE_{11}]\)，不需要除以 \(r\)。即使特征整除 \(r\)，\([I_r]\) 的迹为零，也不妨碍整个同调空间是一维。由\cref{b4:prop:hochschild-low-degrees}，矩阵代数的导子都是内导子；由\cref{b4:thm:hh2-deformation-classification}，其一阶形变均等价于平凡形变。
\end{example}

\subsection{光滑交换代数的 Hochschild 同调}
\label{b4:subsec:1204:02}

\(\operatorname{HH}_1(A)=\Omega_{A/k}\) 提示：在合适的交换代数上，高次同调应当与微分的外幂有关。先看一个可以只用已学消解算完的情形。

\begin{proposition}\label{b4:prop:polynomial-hochschild}
设 \(A=k[x_1,\ldots,x_d]\)，其中 \(k\) 为域，\(d\ge0\)。对每个 \(n\ge0\)，作为 \(A\) 模有
\[
\operatorname{HH}_n(A)\cong\bigwedge_A^n A^d,
\qquad
\operatorname{HH}^n(A)\cong
\bigwedge_A^n\operatorname{Hom}_A(A^d,A).
\]
在同调端可把标准基与 \(\mathrm{d}x_1,\ldots,\mathrm{d}x_d\) 对应；两者在 \(n>d\) 时都为零。
\end{proposition}
\begin{proof}
包络代数为 \(S=k[x_1,\ldots,x_d,y_1,\ldots,y_d]\)，\(A\) 是商 \(S/(x_1-y_1,\ldots,x_d-y_d)\)。这里 \(x_i\) 来自张量积的第一份 \(A\)，\(y_i\) 来自第二份；乘法映射把两份变量送到同一个 \(x_i\)，其核因而由这些差生成。依次令 \(x_i=y_i\) 后仍为多项式环，所以下一个 \(x_j-y_j\) 不为零因子。这是正则列，由\cref{b4:prop:koszul-resolution}，相应 Koszul 复形是 \(A\) 的有限自由 \(S\) 模消解。

例如 \(d=2\) 时，记 \(z_i=x_i-y_i\)，该消解为
\[
0\longrightarrow S\xrightarrow{\mathrm{d}_2}S^2
\xrightarrow{\mathrm{d}_1}S\longrightarrow A\longrightarrow0,
\quad
\mathrm{d}_2(c)=(-z_2c,z_1c),\quad
\mathrm{d}_1(u,v)=z_1u+z_2v.
\]
与 \(A\) 张量后，\(z_1,z_2\) 均为零，所以微分消失，同调依次为 \(A,A^2,A\)，对应于 \(1\)、\(\mathrm{d}x_1,\mathrm{d}x_2\) 及 \(\mathrm{d}x_1\wedge\mathrm{d}x_2\)。一般 \(d\) 时机制相同：张量后各项为 \(\bigwedge_A^n A^d\)，取 \(\operatorname{Hom}_S(-,A)\) 后则为其 \(A\) 线性对偶，两种复形的微分都为零。有限自由模的楔积基将后者识别为 \(\bigwedge_A^n\operatorname{Hom}_A(A^d,A)\)。微分模的基 \(\mathrm{d}x_i\) 则由多项式逐变量求导的泛性质得到。
\end{proof}

要推广到一般交换代数，需要区分光滑与仅仅没有明显奇异方程。下面给出比较定理所需的精确条件。

\begin{definition}\label{b4:def:smooth-algebra-hh}
设 \(k\) 为域，\(A\) 为有限展示的交换含幺 \(k\) 代数。若对每个交换含幺 \(k\) 代数 \(E\)、每个平方零理想 \(J\subseteq E\) 及每个 \(k\) 代数同态 \(A\to E/J\)，都存在提升 \(A\to E\)，则称 \(A\) 在 \(k\) 上\term[guanghua-daishu]{光滑}[smooth]。平方零指 \(J^2=0\)，提升不要求唯一。
\end{definition}

这个约定与第三卷定义14.5.1、14.5.2相同：平方零提升的存在性称为形式光滑，有限展示加形式光滑称为光滑。在形变问题中，用提升性质表述光滑，能够直接看出方程中的误差是否可消去。例如 \(A=k[x]/(x^3)\) 不光滑：取 \(E=k[z]/(z^4)\)、\(J=(z^3)\)，则 \(J^2=0\)，且 \(E/J\cong A\)。恒等映射若能提升，\(x\) 的像必为 \(z+c z^3\)，但其立方是 \(z^3\ne0\)，无法保持 \(x^3=0\)。这在任意特征下都给出一个具体的提升失败。

\ProseParagraphBreak
多项式代数满足此条件，因为分别提升各生成元的像即可。这里要求 \(E\) 交换；它不等同于对所有非交换扩张的提升性质，也不蕴含 \(\operatorname{HH}^2(A)=0\)。例如两个变量的多项式代数在\cref{b4:prop:polynomial-hochschild}中有非零二次上同调。

\ProseParagraphBreak
以下 Hochschild--Kostant--Rosenberg 定理\footnote{科斯坦特（Bertram Kostant，1928-2017），美国数学家，研究李群与表示论；罗森伯格（Alex F. T. W. Rosenberg，1926-2007），德裔美国数学家，研究环论和同调代数。}给出证明概要。Weibel\cite[Theorem 9.4.7]{Weibel1994HomologicalAlgebra}给出更一般的本质有限型版本；这里仅陈述有限展示的情形。

\begin{theorem}[Hochschild--Kostant--Rosenberg 定理]\label{b4:thm:hkr-external}
设 \(k\) 为域，\(A\) 为在 \(k\) 上光滑的有限展示交换代数。则有自然 \(A\) 模同构
\[
\bigwedge_A^n\Omega_{A/k}\longrightarrow\operatorname{HH}_n(A)
\qquad(n\ge0).
\]
它把 \(a_0\cdot\mathrm{d}a_1\wedge\cdots\wedge \mathrm{d}a_n\) 送到
\[
\left[\sum_{\sigma\in S_n}\operatorname{sgn}(\sigma)
a_0\otimes a_{\sigma(1)}\otimes\cdots\otimes a_{\sigma(n)}\right].
\]
\end{theorem}

\begin{proofsketch}
先构造定理中的全局映射，再局部检验它是否为同构。交换代数的 Hochschild 链带有 shuffle 乘法：将两条链的首因子相乘，把其余两串因子按各自原有次序交错排列，并赋予交错置换的符号。该乘法与边界相容，因而下降到同调；这里相乘的是同调链的类，不能使用前面将余链值相乘的杯积公式。最低阶计算为
\[
[1\otimes a]\,[1\otimes b]
=[1\otimes a\otimes b-1\otimes b\otimes a].
\]
这两项分别对应 \(a,b\) 的两种交错次序。特别地，同一个一次链的 shuffle 平方为零，包括特征二的情形；这是交错项在链上配对抵消的结果，不能仅由分次交换律推出。\Cref{b4:prop:hh1-differentials}已使 \(\mathrm{d}a\mapsto[1\otimes a]\) 保持微分模的关系。shuffle 乘法因此诱导自然 \(A\) 线性映射
\(\Psi:\bigwedge_A^*\Omega_{A/k}\to\operatorname{HH}_*(A)\)。反复相乘一次类，便得到各张量因子的带符号置换和，正是定理的公式。这同时证明了公式对微分模关系和外幂关系良定义，没有使用 \(n!\) 可逆。

写 \(A^e=A\otimes_kA\)，以乘法映射的核 \(J\) 定义对角理想。令 \(\Delta a=a\otimes1-1\otimes a\)，则
\[
\Delta(ab)=(a\otimes1)\Delta b+(1\otimes b)\Delta a.
\]
模 \(J^2\) 后，两份 \(b\) 的作用相同，故 \(a\mapsto\Delta a+J^2\) 是 \(k\) 导子，诱导 \(\Omega_{A/k}\to J/J^2\)。其逆把 \(\sum_i a_i\otimes b_i\in J\) 的类送到 \(-\sum_i a_i\cdot\mathrm{d}b_i\)；乘积法则和 \(\sum_i a_ib_i=0\) 保证它零化 \(J^2\)。两次复合在生成元 \(\mathrm{d}a\) 与 \(\Delta a+J^2\) 上为恒等，因而 \(J/J^2\cong\Omega_{A/k}\)。

现在局部计算 \(\operatorname{HH}_n(A)=\operatorname{Tor}^{A^e}_n(A,A)\)。这里使用光滑代数的一项几何结果：在对角上的每个点，\(J\) 在相应的局部环中由正则列生成，其长度等于 \(\Omega_{A/k}\) 的局部秩。这就是光滑有限型态射的对角为正则嵌入的代数表述。从平方零提升条件得到这一局部正则列的证明，见 Weibel\cite[Proposition 9.4.6 and proof of Theorem 9.4.7, pp. 322--323]{Weibel1994HomologicalAlgebra}。

用该正则列的 Koszul 复形消解相应的局部商，再与商张量。正则列作用为零，所以微分全部为零，局部 Tor 代数是 \(J/J^2\) 的外代数。这与多项式计算相同。Koszul 外积和 bar 的 shuffle 乘法相比较时，一次项是刚才的 \(J/J^2\cong\Omega_{A/k}\) 识别，故局部同构就是 \(\Psi\) 的局部化。两个消解的外积比较见\cite[Application 8.7.12, Example 8.7.13, pp. 291--292; Corollary 9.4.4 and proof of Theorem 9.4.7, pp. 322--323]{Weibel1994HomologicalAlgebra}。于是 \(\Psi\) 的核、余核在每个素理想处局部化都为零，因而全局也为零。概要省略的是正则对角的局部证明及 Koszul、bar 外积比较的链映射构造；全局映射本身已由 shuffle 定义，不需另选局部同构来粘合。
\end{proofsketch}

该同构不要求特征零。若另有 \(\operatorname{char}k=0\)，逆映射由下式给出：
\[
[a_0\otimes\cdots\otimes a_n]\longmapsto
\dfrac{1}{n!}a_0\cdot\mathrm{d}a_1\wedge\cdots\wedge \mathrm{d}a_n.
\]
未除以 \(n!\) 的链映射 \(a_0\otimes\cdots\otimes a_n\mapsto a_0\cdot\mathrm{d}a_1\wedge\cdots\wedge\mathrm{d}a_n\) 与正向反对称映射复合后，得到 \(n!\) 倍的恒等映射：每个置换的符号与外幂换序的符号相乘均为正号，共有 \(n!\) 项。因此逆映射的上述写法需要除以 \(n!\)，正向映射的同构性却不需要。例如特征二时，在 \(k[x,y]\) 的二次同调中，\([1\otimes x\otimes y-1\otimes y\otimes x]\) 仍对应非零的 \(\mathrm{d}x\wedge\mathrm{d}y\)，但未除以 \(2!\) 的反向映射把它送到零。上述结论比较的是同调群；整个上链复形的乘法与括号的相容性还需另行讨论。

\subsection{一阶形变的高阶延拓障碍}
\label{b4:subsec:1204:03}

\cref{b4:exm:dual-numbers-first-deformation}的一阶形变确能继续延拓：\(k[\![t]\!][x]/(x^2-t)\) 以 \(1,x\) 为自由基，结合性来自多项式商。但下面的同样小的代数已经有不能通过第二阶的形变。

\begin{example}\label{b4:exm:obstructed-deformation}
设 \(k\) 为任意域，\(A=k\oplus kx\oplus ky\)，其中 \(V=kx\oplus ky\) 满足 \(V^2=0\)，\(1\) 为单位。定义归一化二次上链 \(\phi\)，在 \(V\times V\) 上规定
\[
\phi(x,x)=y,\qquad \phi(y,x)=y,\qquad
\phi(x,y)=\phi(y,y)=0.
\]
当三个自变量属于 \(V\) 时，\(\delta\phi\) 的左右作用项因 \(V^2=0\) 消失，内部项也因相邻乘积为零而消失；某个自变量为 \(1\) 时，归一化使其余两项抵消。由三线性，\(\delta\phi=0\)。故 \(a*b=ab+t\cdot\phi(a,b)\) 确是一阶形变。

\ProseParagraphBreak
然而
\[
(\phi\circ\phi)(x,x,x)
=\phi(y,x)-\phi(x,y)=y\ne0.
\]
对任意二次上链 \(\psi\)，交换性和 \(x^2=0\) 给出
\[
(\delta\psi)(x,x,x)
=x\cdot\psi(x,x)-\psi(x,x)\cdot x=0.
\]
仅有障碍上链在 \((x,x,x)\) 处非零，还不足以断言障碍类非零；关键是所有余边在该处都取零值，而 \(\phi\circ\phi\) 取值为 \(y\)。因此它不可能是余边，\eqref{b4:eq:second-deformation-obstruction}无解。这个一阶形变在任何特征下都不能延拓至模 \(t^3\)。
\end{example}

此例还说明，障碍不只是一句存在性警告。给出乘法表后，先计算\(2\)-余圈，再在一个特定三元组上计算障碍，最后证明所有余边在该三元组上的值都为零，就足以判定无法延拓。其余未选的二阶系数不能改变这一结论。
